\chapter{数据结构}
\bookmark{bookmark:mark-1}{mark-1}
欢迎来到我的大陆。请您放心，我不会在此处引入诸如“数据”的定义、什么是“数据结构三要素”等等内容。但是
同时也请您放心，我会让您彻底搞懂，什么是“数据结构”，“数据结构”究竟是在研究什么。% Latex中，通过\\ 另起一行。

\section{数据结构}
\label{sec:数据结构}
事实上，“数据结构”看起来抽象晦涩的定义，恰恰能够最好地理解“什么是数据结构”。接下来就请你随我一起进入这段拨云见日的旅程。
\begin{definitionbox}
    \term{数据结构}是相互之间存在一种或多种特定“关系”的数据元素的集合，以及定义在这些数据元素上的“操作”的集合。
\end{definitionbox}
我们以一个例子开启这段旅程：假设你是一个游戏程序员，你要编写一个游戏，玩家操控具有四个“角色”组成的小队，如：邪念、银心、
甘尔以及张三这四个人，那么，每个“角色”就都可以被看成是一个数据元素。
\begin{intuitionbox}
\textbf{数据元素}，就是我们要研究和处理的一个个“东西”。这些“东西”可以是一个游戏角色、一件装备、一条聊天消息、
一个整数、甚至一张地图等等。
\end{intuitionbox}
但是，仅仅有了四个数据元素，还不能叫“数据结构”，我们还要解决一个问题：这四个人是什么关系？比如你设计的游戏玩法是将
四人排成一个有前后顺序的队伍，那么四人便会有如下的\uwave{线性关系}：
\begin{bulk}
\begin{center}
邪念 → 银心 → 甘尔 → 张三
\end{center}
\end{bulk}
或者你设计的玩法是一个队长三个队员，则四个人便有如下的\uwave{树状关系}：
\begin{bulk}
\begin{center}
\begin{tikzpicture}
[
    level distance=1.2cm,
    sibling distance=2cm
]
\node[treenode]{邪念}
    child {
        node[treenode]{银心}
    }
    child {
        node[treenode]{甘尔}
    }
    child {
        node[treenode]{张三}
    };
\end{tikzpicture}
\end{center}
\end{bulk}
再假设，四个人之间互相存在好感度，则四个人可能有如下的\uwave{图结构}：
\begin{bulk}
\begin{center}
\begin{tikzpicture}
% 节点
\node[graphnode] (A) {邪念};
\node[graphnode] (B) [right=2cm of A] {银心};
\node[graphnode] (C) [below=1.5cm of A] {甘尔};
\node[graphnode] (D) [below=1.5cm of B] {张三};

% 边 + 权值
\draw (A) -- node[above] {5} (B);
\draw (A) -- node[left] {2} (C);
\draw (A) -- node[right] {8} (D);
\draw (B) -- node[right] {3} (D);
\draw (C) -- node[below] {6} (D);
\end{tikzpicture}
\end{center}
\end{bulk}
请注意，这里我们使用\uwave{无向图}描述，是因为我们假设了A对B的好感度与B对A的好感度相同，但现实中往往并非如此，要描述这种更
现实的情况，就要使用{\color{red}\textbf{后面将会介绍到的}}\uwave{有向图}。

到现在，“数据结构”这个词是不是已经变得具体了？再回看数据结构的定义，我们看到它还有后半句：“操作”。

假设现在我们的游戏队伍：
\begin{bulk}
\begin{center}
邪念 → 银心 → 甘尔 → 张三
\end{center}
\end{bulk}
玩家可能会在游戏中让某些人“加入队伍”（Insert），某些人“退出队伍”（Delete），如卡姐加入：
\begin{bulk}
\begin{center}
邪念 → 银心 → 甘尔 → 张三 → 卡姐
\end{center}
\end{bulk}
邪念退出：
\begin{bulk}
\begin{center}
银心 → 甘尔 → 张三 → 卡姐
\end{center}
\end{bulk}
这里“入队”、“出队”，就都是所谓的“定义在这些数据元素上的操作”。再比如，
新人入队前，你可能需要判断当前队伍是否已满，因此需要“获取队伍人数”（GetSize），这也是“定义在这些数据元素上的操作”。

有了以上的介绍，再回过头来看那个之前觉得它晦涩难懂的数据结构定义，就豁然开朗了。

任何一本《数据结构》的教材，都一定会有“线性表”、“队列”、“树”、“图”等这些内容。为什么我们要重点研究这些数据结构呢？
经过简单的思考不难发现，“线性表”、“队列”、“树”、“图”等，正是我们处理实际需求时，组织数据所需要的最有用的、最常见的结构。
比如，一个有趣的事实是：你的个人计算机中目录结构的组织，其实就是一棵树。如：
\begin{bulk}
\begin{center}
\begin{tikzpicture}
[
    level distance=1.2cm,
    level 1/.style={
        sibling distance=6cm
    },
    level 2/.style={
        sibling distance=1.8cm
    },
    level 3/.style={
        sibling distance=1.5cm
    },
]
\node[treenode]{此电脑}
    child {
        node[treenode]{本地磁盘(C:)}
        child{
            node[treenode]{Program Files}
        }
        child{
            node[treenode]{Windows}
        }
        child{
            node[treenode]{用户}
        }
        child {
            node[treenode]{...}
        }
    }
    child {
        node[treenode]{新加卷(D:)}
        child{
            node[treenode]{电影}
            child {
                node[treenode]{泰坦尼克号}
            }
            child {
                node[treenode]{...}
            }
        }
        child{
            node[treenode]{照片}
        }
        child {
            node[treenode]{...}
        }
    };
\end{tikzpicture}
\end{center}
\end{bulk}
请注意，以上出现的线性结构、树状结构、图结构，都只是逻辑意义上的结构，展示出数据元素之间的逻辑关系，而非数据元素物理存储（如在主存或磁盘中）的位
置关系。这种逻辑上的关系称为数据的\term{逻辑结构}。《数据结构》不仅研究数据的逻辑结构，还要研究如何实现物理存储，这通常通过计算机语言来实现。本书中后面
我们介绍到的每一种数据结构，都会给出其比较常见的、一般的实现方式。如果你感兴趣的话可以想一想，以上出现的数据结构，你如何通过你已有的编程语言知识来实现？

\section{算法与算法效率的度量}
我们已经知道，数据结构研究的是数据的组织方式以及对这些数据的操作。事实上，同样的一批数据，由于组织方式不同，
可能会导致操作效率产生巨大的差异。因此，在实际程序中，我们在考虑数据应该如何组织的同时，也要考虑面对
这些数据，我们应该采用什么步骤和方法去解决问题？例如：
\begin{itemize}[itemsep=0pt,
    parsep=0pt,
    topsep=0pt]
    \item 如何在一组数据中快速找到目标元素？
    \item 如何将一组无序数据排列成有序序列？
    \item 如何在复杂的网络关系中寻找最短路径？
\end{itemize}

对于同一个问题，往往存在多种不同的\uwave{算法}。它们可能都能得到正确结果，但执行过程中的时间消耗、空间占用
可能存在巨大差异。一个算法的效率正是通过其“时间消耗”与“空间占用”来衡量的，正如下面将要介绍的\term{时间复杂度}和
\term{空间复杂度}。

\subsection{时间复杂度}
我们描述一个算法的时间消耗，不能简单地说出类似“算法A需要3秒，算法B需要0.001秒”这样的话。因为代码的运行时间受到
CPU性能等诸多因素的影响。因此，我们需要一种与机器无关的评价方法，这便是\term{时间复杂度}。算法的\term{时间复杂度}，
是指算法执行时间\uline{随问题规模增长而变化}的数量级关系。我们以一个最简单的，在数组中顺序查找某个元素的算法为例来说明：
\begin{cppcode}
for(int i = 0; i < n; i++)
{
    if(a[i] == x)
        return i;
}
\end{cppcode}
这里\(n\)是数组的大小，也即问题规模大小。如何描述该算法的时间复杂度？一个很准确的做法是用该算法的\term{基本操作}的次数来描述。所谓\term{基本操作}，
是指该算法中执行次数最多的、能代表性能瓶颈的操作。该例中，循环执行若干次 i<n, a[i]==x, 以及 i++ 的操作，这两个判断及一个变量自增就是该算法的基本操作。
我们使用\term{\(T(n)\)}来描述算法的基本操作次数，显然，该例中，最坏情况下查找到最后一个，算法才会返回，因此最坏情况下：
\[
T(n)=n
\]
而平均情况下：
\[
T(n) = \frac{1}{n}\cdot 1 + \frac{1}{n}\cdot 2 + ... + \frac{1}{n}\cdot n = \frac{n+1}{2}
\]
这表示若输入规模\(n\)增加一倍，则算法在最坏情况下执行次数也会增加一倍，而平均情况下，执行次数会增加\(\frac{n+1}{2}\)。这便是
时间复杂度的“大T表示法”。

而实际的学术或工程中最最常用的，也相信你一定已经见过了的，并非大T表示法，而是\term{“大O表示法”}：

考虑某算法的时间复杂度为
\[
T(n) = 3n^{2} + 5n + 10
\]
当\(n\)很大时，该算法的时间代价几乎完全由\(3n^2\)这一项决定，而后面两项可以忽略不计。而3是常数，不随问题规模\(n\)的变化而变化。因此，使用
大O表示法，我们说该算法的时间复杂度是\(O(n^2)\)。若某算法的时间复杂度不随问题规模\(n\)的变化而变化，则我们说该算法的时间复杂度为\(O(1)\)。
对于大O表示法严谨的数学定义，感兴趣的同学请自行参考相关资料。
根据我们已有的数学知识，不难得出，在问题规模\(n\)很大的情况下：
\begin{bulk}
    \(O(1)\) 优于 \(O(logn)\) 优于 \(O(n)\) 优于 \(O(nlogn)\) 优于 \(O(n^2)\) 优于 \(O(n^3)\) 优于 \(O(a^n)\)
\end{bulk}
\begin{insightbox}
这里还有一个拉扯感比较强的问题：所谓“问题规模很大”，“很大”的边界在哪？比如某个问题我们确定问题规模\(n\)一定满足\(n<=100\)，那么，
\(T(n) = n^3\), 即\(O(n) = n^3\)的算法，是不是反而优于\(T(n) = 101n^2\), 即\(O(n) = n^2\)的算法？答案是：是的。这个问题也说明了，大O复杂度
并不是算法实际运行效率的绝对比较。但是在实际的工程中，性能瓶颈通常是由大规模（n很大）的数据部分产生的，这也是实际中都在使用大O复杂度的原因。
\end{insightbox}

\subsection{空间复杂度}
某算法的\term{空间复杂度}是指实现该算法所需额外存储空间\uline{随问题规模增长而变化}的数量级关系。同时间复杂度一样，绝大多数情况下我们关注的是算法
所需额外存储空间随问题规模\(n\)的增长趋势，因此我们同样使用大O表示法来描述空间复杂度。

请注意，计算空间复杂度时我们只考虑处理问题而额外引入的空间，并不计算输入数据本身的空间。比如以下数组求和的例子，我们只引入了存储变量 i 和 result 的
两个额外 int 大小的空间，该空间大小与问题的输入规模\(n\)无关，因此该算法的空间复杂度为O(1)，而数组a[]是问题本身的数据，其存储空间并不算入空间复杂度：
\begin{cppcode}
int sum(int a[], int n)
{
    int result = 0;
    for(int i=0; i<n; i++)
    {
        result += a[i];
    }
    return result;
}
\end{cppcode}

\begin{notebox}{提示}
    分析空间复杂度很简单，只需要看程序额外使用了多少内存。如果你对时间复杂度没有“感觉”，笔者在这里给出一个建立感觉的提示。如前所述，时间复杂度要看程序
    执行次数最多、最能代表瓶颈的操作。通常来说，执行次数最多的循环就是程序的瓶颈。如：
    \begin{cppcode}
        void f(int n) { // 模拟问题的输入规模为n
            int a = 10;
            int b = a + 2;
            // 该循环约执行 n 次，复杂度为n量级
            for(int i=0; i<n; i++) {
                // ...
            }
            // 该循环约执行 n+(n-1)+...+2+1 次，复杂度为@\(n^2\)@量级
            for(int i=0; i<n; i++) {
                for(int j=i; j<n; j++) {
                    // ...
                }
            }
            // 该循环约执行 nlogn次，复杂度为nlogn量级
            for(int i=0; i<n; i++) {
                for(int j=1; j<n; j*=2) {
                    // ...
                }
            }
        }
    \end{cppcode}
    以上代码中，显然
    \begin{cppcode}
        // 该循环约执行 n+(n-1)+...+2+1 次，复杂度为@\(n^2\)@量级
        for(int i=0; i<n; i++) {
            for(int j=i; j<n; j++) {
                // ...
            }
        }
    \end{cppcode}
    为程序的性能瓶颈。因此我们说算法f的时间复杂度是\(O(n^2)\)。
\end{notebox}

\section{\color{red}\textbf{标题还没想好}}
通过以上的介绍，我们就了解了《数据结构》或者叫《数据结构与算法》这个学科，就是研究在不同的实际需求下，数据如何存储、组织，以及高效操作的学科。

任何一本数据结构的教材都会重点介绍“线性表”、“队列”、“树”、“图”等数据结构，这是因为这些结构并不是枯燥的理论模型，而是计算机世界中最常见、最高效的组织方式。
当你打开一个游戏背包，当你浏览一棵文件目录树，当你在地图中寻找一条最短路径，当你在社交网络中查看好友关系，你看到的表象虽然不同，但它们背后都隐藏着数据结构的思想。

笔者认为，《数据结构》这个学科次在“掌握知识”，而重在“培养能力”。学习《数据结构》，真正重要的是，在学习的过程中不断思考、体会和理解这些数据结构与算法背后的思想及其
代码实现原理，训练自己灵活的思维和解决问题的能力。这也是在本书中我将一遍又一遍地强调的内容。

那接下来，就我们将一起走进这段探索之旅，逐一认识这些经典结构，理解它们为什么诞生，以及它们如何帮助我们更加高效地解决实际问题。在此之前，如果你是一个喜欢思考问题的人，
不妨看看下面的“思考”。
说不定天才的你能够自己探索出最优解（甚至想出了比目前已知的最优解更优的解法，顺便发表一篇论文什么的）。以下的这些问题呢，也会随着本书后续的开展，
而不断地展开讨论，相信你每次回看它们的时候，都会有全新的不同理解。
\begin{thinkbox}{想一想}
{\color{red}\textbf{这里给出一些思考吧，学习本书之前，可以自己瞎想一下，学完本书之后，都给出一个很优的答案，尽量做到每种结构一个例子。}}
\end{thinkbox}


\begin{cppcodeNum}[label={code:test}, codetitle={测试代码}]

int main() {
    int a = 10;
    int b = 20;

    int c = a + b;

    return 0;
}

\end{cppcodeNum}

带页码：如\figref{fig:顺序存储的线性表}所示。如\coderef{code:test}所示。
不带页码：如图\ref{fig:顺序存储的线性表}所示。如代码\ref{code:test}所示。